%! AP Statistics — Unit 1, Lesson 1.2 — Homework KEY
\documentclass[10pt]{article}
\usepackage{apstats-article}
\usepackage{apstats-key}

%=============================================================================
\begin{document}

\pageheader{Unit 1, Lesson 1.2}{Homework --- Answer Key}
\begin{tabularx}{\linewidth}{X X X}
Name:\;\ans{Key} & Date:\; & Period:\;
\end{tabularx}

\vspace{0.1in}

% ════════════════════════════════════════════════════════════════════════════
% PART A
% ════════════════════════════════════════════════════════════════════════════
\begin{blurbbox}[Part A \quad NBA Player Dataset]{navylight}
\small

\renewcommand{\arraystretch}{1.5}
\begin{tabularx}{\linewidth}{@{} >{\bfseries}p{0.16\linewidth}
                                    C{0.13\linewidth} C{0.13\linewidth}
                                    C{0.16\linewidth} C{0.13\linewidth}
                                    C{0.14\linewidth} @{}}
\rowcolor{sky}
\textbf{Name} & \textbf{Team} & \textbf{Pts/game} &
\textbf{Position} & \textbf{Games} & \textbf{Jersey \#} \\
\midrule
Player A & Lakers  & 27.4 & Forward & 68 & 23 \\
Player B & Celtics & 19.8 & Guard   & 72 & 11 \\
Player C & Warriors& 31.2 & Guard   & 65 & 30 \\
Player D & Heat    & 14.5 & Center  & 79 & 3  \\
Player E & Bucks   & 23.7 & Forward & 71 & 34 \\
\end{tabularx}

\vspace{10pt}

\renewcommand{\arraystretch}{2.0}
\begin{tabularx}{\linewidth}{@{} >{\bfseries}p{0.2\linewidth}
                                    C{0.14\linewidth} X @{}}
\rowcolor{sky}
\textbf{Variable} & \textbf{Type} & \textbf{Justification} \\
\midrule
Name         & \ans{C}  & \ans{A label identifying the individual; no arithmetic applies} \\
Team         & \ans{C}  & \ans{Places player into a named category; averaging team names is meaningless} \\
Pts/game     & \ans{QC} & \ans{Measured scoring rate; can take any non-negative decimal value} \\
Position     & \ans{C}  & \ans{Labels a role category (Guard, Forward, Center)} \\
Games played & \ans{QD} & \ans{Count of whole games; gaps between values (can't play 68.5 games)} \\
Jersey \#    & \ans{C}  & \ans{Numeric label; averaging jersey numbers gives no meaningful result} \\
\end{tabularx}

\vspace{8pt}
\textbf{Which variable is most likely to cause confusion?}\\[4pt]
\ans{Jersey number --- it looks quantitative because it is a number, but it is a label.
Averaging jersey numbers (e.g., mean jersey = 20.2) produces a value that means nothing
about the players. This is the key ``trap'' variable of this lesson.}

\end{blurbbox}

\begin{teachernote}
Pts/game: a small number of students may say QD because scoring is tracked in
half-points in some contexts. Award credit for QC (the standard classification) and
accept QD with a specific justification about how points are recorded in the source data.
\end{teachernote}

\vspace{0.14in}

% ════════════════════════════════════════════════════════════════════════════
% PART B
% ════════════════════════════════════════════════════════════════════════════
\begin{blurbbox}[Part B \quad AP-Style Multiple Choice]{greenacc}
\small

\textbf{Question 1.}\enspace Grade level (9, 10, 11, 12).\\[6pt]
\begin{enumerate}[label=\textbf{(\Alph*)}, leftmargin=2em, itemsep=3pt]
  \item Grade level is quantitative continuous.
  \item Grade level is quantitative discrete.
  \item \textbf{\textcolor{keyred}{Grade level is categorical.}} \textcolor{keyred}{$\leftarrow$ \textbf{correct}}
  \item Cannot be classified.
\end{enumerate}
\ans{Explanation: Grade level uses numbers to label groups (freshmen = 9, sophomore = 10, etc.).
Averaging grade levels (e.g., mean grade = 10.4) tells us nothing meaningful about student
membership in a class year. The numbers are identifiers, not measured amounts.}\\[4pt]

\begin{teachernote}
Choice B is the most common wrong answer. Students see whole numbers and default to QD.
The key question: ``Does the average of these values mean something?'' An average grade
level of 10.3 does not correspond to a real grade level. Push students to articulate why.
\end{teachernote}

\vspace{10pt}
\textbf{Question 2.}\enspace Type of precipitation + total amount (inches).\\[6pt]
\begin{enumerate}[label=\textbf{(\Alph*)}, leftmargin=2em, itemsep=3pt]
  \item Both categorical.
  \item \textbf{\textcolor{keyred}{Type = categorical; amount = quantitative continuous.}}
        \textcolor{keyred}{$\leftarrow$ \textbf{correct}}
  \item Type = QD; amount = QC.
  \item Both quantitative.
\end{enumerate}
\ans{Explanation: Type of precipitation (rain, snow, sleet, none) places each day into a
named category --- no arithmetic applies. Total precipitation is a measured amount in inches
that can take any non-negative real value, making it quantitative continuous.}

\vspace{10pt}
\textbf{Question 3.}\enspace Pain level (0--10) + days since surgery.\\[6pt]
\begin{enumerate}[label=\textbf{(\Alph*)}, leftmargin=2em, itemsep=3pt]
  \item Both QD.
  \item Pain = categorical; days = QD.
  \item \textbf{\textcolor{keyred}{Pain = debatable (C or QD); days = QD.}}
        \textcolor{keyred}{$\leftarrow$ \textbf{best answer}}
  \item Both continuous.
\end{enumerate}
\ans{Explanation: Days since surgery is a count of whole days --- clearly QD. Pain level is
an ordinal rating scale; whether equal spacing can be assumed is debatable. Most AP Statistics
resources treat 0--10 pain scales as categorical (ordinal), but some analysts treat them as
QD. The classmate is partially correct about days since surgery but oversimplifies pain level.}

\end{blurbbox}

\vspace{0.14in}

\begin{extensionbox}
\small
\textbf{Optional Extension:}\\[4pt]
\ans{Answers vary by dataset chosen. Award credit for any six variables correctly classified with
justification. Look for: (1) at least one non-obvious categorical variable (e.g., FIPS code, race
category), (2) correct identification of discrete vs.\ continuous among quantitative variables
(e.g., count of deaths = discrete; age in years as reported = debatable; percent = continuous).}\\[4pt]
\textcolor{slate}{\scriptsize Source: student-chosen from CDC or Census website.}
\end{extensionbox}

\end{document}
